\subsection{Miscellaneous Formulas}

\begin{lemma}
    \label{cartan_calc:lem:relation_between_Lie_derivative_and_connection}
    Let $M$ be a Riemannian manifold, $\omega \in \Omega^k(M)$ a $k$-form and $X \in \mathscr{X}(M)$ a vector field. Then
    \[
    \mathcal{L}_X \omega = \nabla_X \omega + \sum_{i=1}^{k} \omega(-,\ldots,\nabla X,\ldots,-),
    \] 
    where $\nabla X$ enters in the $i$-th slot in the $i$-th summand, which is a covariant tensor field of degree $k$ defined by 
    \[
    \omega(-,\ldots,\nabla X,\ldots,-)(Y_1,\ldots,Y_k) := \omega(Y_1,\ldots,\nabla_{Y_i} X,\ldots,Y_k).
    \] 
\end{lemma}

\begin{proof}
    Let $e_1,\ldots,e_n$ be local orthonormal frame and let $\epsilon^1,\ldots,\epsilon^n$ be the corresponding dual frame. Using Cartan's magic formula (Proposition \ref{cartan_cal:prop:magic_formula}) and the formula for the exterior derivative from Lemma \ref{cartan_calc:lem:formula_for_d_and_codifferential} we compute 
    \begin{align*}
    \mathcal{L}_X\omega & = d\iota_X\omega + \iota_X d \omega \\
                        & = \sum_{i=1}^{n}  \epsilon^i \wedge \nabla_i \iota_X\omega + \iota_X \left(\sum_{i=1}^{n} \epsilon^i \wedge \nabla_i \omega\right) \\
                        & = \sum_{i=1}^{n}  \epsilon^i \wedge \nabla_i \iota_X\omega + \sum_{i=1}^{n} \iota_X \epsilon^i \wedge \nabla_i \omega - \epsilon^i \wedge \iota_X \nabla_i \omega 
    \end{align*}
    In the last line we used the Leibniz rule for the interior product, see Theorem \ref{thm:cartan_calculus}. Moreover 
    \[
    \iota_X \nabla_i \omega = \nabla_i(\omega(X,-)) - \omega(\nabla_iX,-).
    \]
    Plugging this into the formula above we see that the first sum cancles and we obtain 
    \[
    \mathcal{L}_X\omega = \sum_{i=1}^{n} \iota_X \epsilon^i \wedge \nabla_i \omega + \epsilon^i \wedge \omega(\nabla_iX,-)
    \]
    To compute the first summand write $X = \sum_{j=1}^{n}X^je_j$. Then 
    \[
    \sum_{i=1}^{n} \iota_X \epsilon^i \wedge \nabla_i \omega = \sum_{i=1}^{n} X^i \nabla_i \omega = \nabla_X \omega.
    \]
    For the second summand let $e_{j_1},\ldots,e_{j_k}$ be any (not necessarily ordered) $k$-tuple of distinct basis vectors. If $i = j_l$ for some $1 \leq l \leq k$ we get 
    \begin{align*}
    (\epsilon^i \wedge \omega(\nabla_iX,-))(e_{j_1},\ldots,e_{j_k}) & = (-1)^{l-1}\omega(\nabla_iX,e_{j_1},\ldots,e_{j_{l-1}},e_{j_{l+1}},\ldots,e_{j_n}) \\ 
    &= \omega(e_{j_1},\ldots,e_{j_{l-1}},\nabla_iX,e_{j_{l+1}},\ldots,e_{j_n}) 
    \end{align*}
    On the other hand, if $i$ is distinct from any of the $j_l$, this term evaluates to $0$. Therefore
    \begin{align*}
        \left(\sum_{i=1}^{n} \epsilon^i \wedge \omega(\nabla_iX,-)\right)(e_{j_1},\ldots,e_{j_k}) &= \sum_{l=1}^k \omega(e_{j_1},\ldots,e_{j_{l-1}},\nabla_{j_l}X,e_{j_{l+1}},\ldots,e_{j_n}) \\
        &= \left(\sum_{i=1}^{k} \omega(-,\ldots,\nabla X,\ldots,-)\right)(e_{j_1},\ldots,e_{j_k})
    \end{align*}
    % \begin{align*}
    %     (\epsilon^i \wedge \omega(\nabla_iX,-))(e_1,\ldots,e_n) & = (-1)^{i-1} \omega(\nabla_iX,e_1,\ldots,e_{i-1},e_{i+1},\ldots,e_n) \\
    %     & = \omega(e_1,\ldots,e_{i-1},\nabla_iX,e_{i+1},\ldots,e_n) \\
    %     & = \omega(-,\ldots,\nabla X,\ldots,-)(e_1,\ldots,e_k) 
    % \end{align*}
    Combining these observations the lemma follows.
\end{proof}

\begin{lemma}
    Let $M$ be a Riemannian manifold and $X \in \mathscr{X}(M)$ a vector field. Then 
    \[
    \mathrm{div}(X) = \mathrm{tr}(\nabla X).
    \]
    In particular if $e_1,\ldots,e_n$ is a local orthonormal frame we obtain 
    \[
    \mathrm{div}(X) = \sum_{i=1}^{n} \langle \nabla_{e_i}X, e_i \rangle.
    \]
\end{lemma}

\begin{proof}
    The divergence of a vector field $X$ is defined by 
    \[
    d\iota_X \mathrm{vol} = \mathrm{div}(X) \mathrm{vol},
    \]
    where $\mathrm{vol}$ is the Riemannian volume form (if $M$ is not oriented $\mathrm{vol}$ can only be defined locally, but the divergence will be independent of the chosen local orientation, hence still globally defined). Using Cartan's magic formula and the fact that $\mathrm{vol}$ is closed we obtain from Lemma \ref{cartan_calc:lem:relation_between_Lie_derivative_and_connection}
    \begin{align*}
        \mathrm{div}(X) \mathrm{vol} & = \mathcal{L}_X \mathrm{vol} \\
    & = \nabla_X \mathrm{vol} + \sum_{i=1}^{k} \mathrm{vol}(-,\ldots,\nabla X,\ldots,-) \\
    & = \sum_{i=1}^{k} \mathrm{vol}(-,\ldots,\nabla X,\ldots,-)
    \end{align*}
    The last equality follows from the fact that the volume form is parallel. Let $e_1,\ldots,e_n$ be an oriented orthonormal frame and $\epsilon^1,\ldots,\epsilon^n$ the corresponding dual frame (again if $M$ is not oriented we only work with a local orientation). Write  
    \[
    \nabla X = \sum_{i,j=1}^{n} Y_i^j \epsilon^i \otimes e_j.
    \]
    Then 
    \begin{align*}
        \mathrm{vol}(-,\ldots,\nabla X,\ldots,-)(e_1,\ldots,e_n) & = \mathrm{vol}(e_1,\ldots,\nabla_{e_i}X,\ldots,e_n)  \\
        & = \sum_{j=1}^{n} Y_i^j \mathrm{vol}(e_1,\ldots,e_j,\ldots,e_n)  \\
        &= Y_i^i \mathrm{vol}(e_1,\ldots,e_i,\ldots,e_n) = Y_i^i
    \end{align*}
    and thus 
    \[
    \mathrm{vol}(-,\ldots,\nabla X,\ldots,-) = Y_i^i \mathrm{vol}.
    \]
    By definition of the trace $\mathrm{tr}(\nabla X) = \sum_{i=1}^{n} Y_i^i$, from which the first formula in the lemma now follows. The second formula follows directly from the definition of the trace.
\end{proof}

\begin{remark}
    For a smooth function $f \colon M \to \R$ there are two ways to define the Laplacian of $f$, namely as $\mathrm{div}(\mathrm{grad}(f))$ or as the connection Laplacian $\mathrm{tr}_g(\nabla^2 f)$, where we take the trace with respect to the metric $g$ on $M$. By the above Lemma the two agree, using the fact that $\mathrm{grad}(f) = (\nabla f)^\sharp$ and that the covariant derivative commutes with the musical isomorphisms:
    \[
    \mathrm{div}(\mathrm{grad}(f)) = \mathrm{tr}(\nabla\mathrm{grad}(f)) = \mathrm{tr}((\nabla^2f)^\sharp) = \mathrm{tr}_g(\nabla^2 f)
    \]
    The last equality is the definition of the trace operator $\mathrm{tr}_g$.
\end{remark}